\section{Implementación con OpenMP}

\subsection{La directiva}

El pragma que se aplica al triple bucle, en \texttt{modules/tiempo\_mult\_openmp.c}, es:

\begin{verbatim}
#pragma omp parallel for collapse(2)
for (int i = 0; i < N; i++) {
    for (int j = 0; j < N; j++) {
        for (int k = 0; k < N; k++) {
            C[(size_t)i * N + j] +=
                A[(size_t)i * N + k] *
                B[(size_t)k * N + j];
        }
    }
}
\end{verbatim}

\subsection{Justificación a partir de los ejemplos de la especificación}

Antes de escribir el código se revisaron los ejemplos del documento oficial \textit{OpenMP Application Programming Interface Examples v6.0.1} (\texttt{github.com/OpenMP/Examples}), escogiendo las directivas entendiendo el porqué y no al azar. Del índice, tres ejemplos aplican al patrón del kernel:

\begin{center}
\begin{tabular}{|l|p{5.5cm}|p{6.5cm}|}
\hline
\textbf{Ejemplo} & \textbf{Patrón} & \textbf{Veredicto} \\
\hline
\textbf{3.1} \texttt{ploop.1.c} &
\texttt{\#pragma omp parallel for} sobre un solo loop &
Solo paraleliza el loop externo $i$. Con $N=4000$ son 4000 bloques, mal repartidos entre 12 núcleos. \\
\hline
\textbf{3.8} \texttt{collapse.1.c} &
\texttt{\#pragma omp for collapse(2)} sobre dos loops anidados &
\textbf{Elegido.} Une $i$ y $j$ en un único espacio de $N^2$ iteraciones (16 millones para $N=4000$) con buen reparto entre hilos. \\
\hline
\textbf{5.3.5} \texttt{task\_dep.5.c} &
\texttt{\#pragma omp task depend(in/inout)} sobre bloques &
Descartado. Asume un kernel bloqueado y dependencias explícitas entre bloques; más complejo de lo que el caso requiere. \\
\hline
\end{tabular}
\end{center}

\subsection{Por qué cada elemento}

\begin{itemize}
    \item \textbf{\texttt{parallel for} (combinada)}: crea el equipo de hilos y reparte el trabajo del loop en una sola directiva. El ejemplo 3.1 del PDF muestra las dos formas equivalentes (combinada y separada) y la combinada es más limpia.
    \item \textbf{\texttt{collapse(2)}}: es la clave de la elección. Colapsar dos loops anidados significa que el \textit{runtime} los trata como un único espacio de iteraciones de tamaño $N \times N$ y reparte ese espacio entre los hilos. Sin \texttt{collapse} solo se paralelizaría $i$, y con $N=2000$ eso daría 2000 bloques para 12 núcleos: un reparto muy desigual. El loop $k$ queda secuencial dentro de cada par $(i,j)$, que es lo correcto, porque la acumulación del producto punto debe permanecer en un solo hilo.
    \item \textbf{No se usa \texttt{private(k)} explícito}: en C99 las variables de control de un \texttt{for} ya son implícitamente privadas dentro de un loop construct. Declararlo además \emph{falla la compilación}, porque $k$ está declarada dentro del \texttt{for} y no es visible en el punto donde aparece la directiva.
\end{itemize}

\subsection{Por qué no se necesitan otras cláusulas}

Este es el punto donde el conocimiento del algoritmo garantiza la corrección:

\begin{itemize}
    \item $A$, $B$ y $C$ quedan \textbf{shared} por defecto, que es lo correcto: $A$ y $B$ solo se leen, y cada hilo solo escribe en $C[i \cdot N + j]$ para el par $(i,j)$ concreto que le tocó en el espacio colapsado. No hay race condition sobre $C$.
    \item \textbf{No hace falta \texttt{reduction(+:)}}. Una \textit{reduction} se usa cuando varios hilos acumulan sobre la \emph{misma} variable. Aquí cada $C[i][j]$ lo acumula íntegramente un solo hilo y los hilos distintos escriben en celdas distintas.
    \item \textbf{No hace falta \texttt{critical} ni \texttt{atomic}}: no existe región crítica.
    \item \textbf{No hace falta \texttt{barrier}} explícita: la región paralela ya tiene una barrera implícita al final.
\end{itemize}

\subsection{Configuración del número de hilos}

El número de hilos se puede fijar en tres lugares, por orden de precedencia: la cláusula \texttt{num\_threads()} de la directiva, la llamada a \texttt{omp\_set\_num\_threads()} y la variable de ambiente \texttt{OMP\_NUM\_THREADS}. Aquí se usa la variable de ambiente, que es la forma más portable y además permite barrer el número de hilos \emph{sin recompilar} el binario.

\subsection{El flag de compilación no es opcional}

Un detalle que cuesta caro si se pasa por alto: \textbf{las directivas se ignoran silenciosamente si no se compila con \texttt{-fopenmp}}. El programa compilaría sin errores, se ejecutaría, y daría exactamente el mismo resultado secuencial.

Además, el flag de optimización debe ser \textbf{el mismo que usan las demás estrategias} de la comparación. La primera versión de este trabajo compiló el binario de OpenMP con \texttt{-O0} mientras el resto usaba \texttt{-O2}; esa diferencia hacía que el speedup de OpenMP pareciera artificialmente malo (entre 0.71$\times$ y 1.48$\times$) porque la comparación mezclaba dos efectos distintos: el paralelismo y la ausencia de optimización del compilador. Corregido a \texttt{-fopenmp -O2}, el $T=1$ de OpenMP da 76.02\,s contra los 75.65\,s del secuencial de referencia (0.5\% de diferencia), lo que confirma que la base de comparación quedó justa.
